\section{Introdução}\label{sec:introducao}

A madeira acompanha a história da construção desde os primeiros sistemas estruturais e permanece atual em razão de sua ecoeficiência, do baixo consumo de energia em sua produção e da elevada relação resistência-peso \parencite{FRAGA2025}.
Em edificações, a energia incorporada de estruturas de madeira é, em média, 28\% e 47\% inferior à de estruturas equivalentes em concreto e aço, respectivamente \parencite{Schenk2022}.
Sua associação ao concreto e ao aço permite desenvolver sistemas estruturais mistos que combinam as propriedades dos materiais, ampliando as possibilidades de aplicação da madeira e de melhoria do desempenho estrutural \parencite{PASTORI2022}.
Diante das metas globais de descarbonização da construção civil \parencite{Li2025}, a madeira ressurge como alternativa estrutural promissora para pontes de pequeno e médio porte, particularmente em estradas vicinais, nas quais os vãos são reduzidos e a demanda por soluções econômicas e sustentáveis é elevada \parencite{Milani2020}.
Em muitos desses contextos, soluções em concreto e aço esbarram na indisponibilidade local de produção e nas dificuldades logísticas de transporte desses materiais até o canteiro, o que torna a madeira, frequentemente disponível na própria região, a alternativa estruturalmente mais viável.

Nesse cenário, a otimização estrutural computacional consolida-se como instrumento para conciliar desempenho e metas ambientais, já tendo sido empregada em diversas áreas de projeto em engenharia. \textcite{Moraes2022}, por exemplo, empregou otimização metaheurística na redução das emissões de CO\textsubscript{2} associadas ao dimensionamento de vigas pré-moldadas e protendidas de concreto, ao passo que \textcite{FRAGA2025} aplicou conceitos de otimização ao projeto de elementos de madeira em treliças do tipo Fink.

Historicamente, o processo de projeto estrutural evoluiu de cálculos manuais e procedimentos semiautomatizados de desenho assistido por computador para abordagens que integram modelagem paramétrica, análise estrutural automatizada e técnicas de otimização orientadas a dados.
Esse paradigma já vem sendo aplicado com sucesso ao projeto de pontes: acoplando modelagem paramétrica tridimensional a rotinas de otimização global e local, \textcite{deMoraes2025} automatizaram o dimensionamento de vigas de concreto protendido para pontes rodoviárias.
Em escala industrial, \textcite{Rajic2026} desenvolveram um \textit{toolbox} paramétrico que automatiza a geração de modelos de elementos finitos e as verificações normativas de pontes de concreto protendido, aplicado com sucesso em projetos de infraestrutura de grande porte.
Além disso, plataformas dessa natureza têm sido empregadas para avaliação e diagnóstico estrutural de pontes \parencite{Garbowski2023, Alseid2025, Miluccio2026}.

Essa literatura concentra-se, contudo, predominantemente em pontes de concreto.
Frameworks paramétricos para estruturas de madeira, como o de \textcite{Jiang2023}, têm aplicação restrita a componentes de edificações.
Sua extensão a pontes de madeira roliça ainda é pouco explorada, especialmente diante da variabilidade dimensional das peças e de sua influência no comportamento estrutural.
Os estudos disponíveis sobre pontes de madeira concentram-se majoritariamente em levantamentos de campo e tabelas de pré-dimensionamento baseadas em verificações manuais, como o de \textcite{Criado2024}, ou em ensaios dinâmicos e modelagem numérica de pontes existentes, como o de \textcite{BERGENUDD2023115429}, sem incorporar modelagem paramétrica, otimização estrutural ou tratamento de robustez.
Uma revisão sistemática recente sobre a otimização estrutural de componentes de madeira aponta, ainda, que a indisponibilidade de códigos abertos é um dos principais entraves à transferência dessa pesquisa para aplicações práticas \parencite{Belz2025}.
A literatura sobre otimização de pontes permanece, assim, dominada por outros materiais, e para pontes rodoviárias de madeira roliça permanece limitada a investigação integrada entre geometria, propriedades resistentes, variabilidade dimensional, eficiência material e estados limites governantes.
É essa lacuna que o framework paramétrico de otimização robusta proposto neste trabalho busca preencher, com a plataforma computacional de acesso livre servindo como instrumento de implementação e transferência dos resultados.
Nesse paradigma, o dimensionamento deixa de ser um procedimento linear e passa a ser tratado como um processo iterativo e integrado, no qual o modelo geométrico, a análise dos esforços e a decisão de projeto são acoplados dentro de uma única rotina de otimização.
Essa transição permite reduzir o tempo de projeto, ampliar a qualidade das soluções e explorar de forma sistemática o espaço de alternativas viáveis.

Nesse contexto, o dimensionamento de pontes de madeira pode ser formulado como um problema de otimização, no qual objetivos conflitantes devem ser considerados simultaneamente, como a minimização do consumo de material, o atendimento aos estados limites último e de serviço e o desempenho estrutural global.
A otimização multiobjetivo baseada em metaheurísticas surge como ferramenta adequada para lidar com essa complexidade, pois possibilita a obtenção de um conjunto de soluções eficientes representadas por uma fronteira de Pareto \parencite{DEB2009}, oferecendo ao projetista maior flexibilidade na tomada de decisão e favorecendo o uso racional dos recursos.

Entretanto, a otimização determinística conduz, com frequência, a soluções localizadas exatamente sobre a fronteira das restrições, nas quais pequenas variações nas variáveis de projeto, decorrentes de tolerâncias de fabricação, variabilidade dimensional da madeira roliça ou imprecisões de montagem, podem levar à violação dos critérios de segurança.
Esse aspecto é particularmente crítico em madeira roliça, cujo diâmetro apresenta variabilidade natural significativa e cuja capacidade resistente está diretamente condicionada a essa variabilidade dimensional \parencite{Morgado2017}.
Assim, torna-se necessário incorporar o conceito de robustez ao processo de otimização, de modo que as soluções obtidas permaneçam seguras mesmo diante de desvios em torno de seus valores nominais.

\subsection{Objetivos e contribuições}\label{sec:objetivos_contribuicoes}

Este trabalho propõe um framework paramétrico de otimização robusta multiobjetivo para o dimensionamento de pontes de madeira roliça, no qual a modelagem paramétrica dos elementos estruturais é acoplada a uma rotina de busca evolutiva, em conformidade com os critérios da \textcite{NBR7190-1:2022}.
Diferentemente de abordagens que dimensionam o tabuleiro e a longarina em etapas sequenciais e independentes, a formulação aqui adotada acopla as duas peças em uma única avaliação, detalhada na Seção~\ref{sec:metodologia}.

Os objetivos específicos do trabalho são os seguintes.
\begin{enumerate}[label=(\roman*)]
    \item Formular o dimensionamento dos elementos estruturais de pontes de madeira como um problema de otimização multiobjetivo, definindo funções objetivo e restrições compatíveis com os critérios de segurança e desempenho da norma.
    \item Implementar o algoritmo NSGA-II acoplado ao modelo estrutural paramétrico, assegurando a convergência e a diversidade das soluções.
    \item Incorporar a análise de robustez ao processo de busca, avaliando o efeito de um desvio de 5\% aplicado ao diâmetro nominal da longarina sobre as restrições.
    \item Investigar a matriz paramétrica de pontes de estrada vicinal de pista simples, cobrindo quatro vãos teóricos e as cinco classes de resistência de madeiras da norma, de modo a identificar as tendências de consumo de material, desempenho estrutural e estados limites governantes ao longo do domínio investigado.
\end{enumerate}

As principais contribuições deste trabalho podem ser separadas em três frentes: metodológica, estrutural e aplicada. Até onde se tem conhecimento, a contribuição metodológica é a primeira abordagem a integrar modelagem paramétrica, otimização multiobjetivo e análise de robustez no dimensionamento de pontes de madeira roliça. A contribuição estrutural vem da investigação da matriz paramétrica, que identifica as relações entre vão, classe de resistência, variabilidade dimensional, consumo de material e estados limites governantes. A contribuição aplicada resulta da matriz de casos avaliada, que produz ábacos de anteprojeto e uma equação fechada de pré-dimensionamento, disponibilizados por meio de uma plataforma computacional de acesso livre que implementa e transfere esses resultados ao projetista.
